\documentclass[11pt]{article}
\usepackage{amsmath}
\usepackage{amssymb}
\usepackage{geometry}
\geometry{margin=1in}
\title{GRPO Training and Reward Flow for Tesserae (RLVR)}
\author{}
\date{}

\begin{document}
\maketitle

\section{Scope}
This document summarizes how RLVR computes reward and applies GRPO updates when training on the Tesserae dataset. It covers the data pipeline, reward computation, token-level reward tensor, GRPO advantage construction, and the forward/backward flow for multimodal (image + text) inputs. It also explains why the vision tower is updated even though the advantage is applied only to response tokens.

\section{Data and Prompt Pipeline}
\begin{itemize}
  \item Dataset loader: \texttt{rlvr/verl/utils/dataset.py} (\texttt{RLHFDataset}).
  \item Training config: \texttt{rlvr/examples/config.yaml}.
  \item Prompt formatting: \texttt{rlvr/examples/format_prompt/math.jinja}.
\end{itemize}

\subsection{Keys and fields}
The Tesserae training config uses:
\begin{itemize}
  \item \texttt{prompt	extunderscore key: problem}
  \item \texttt{answer	extunderscore key: answer}
  \item \texttt{image	extunderscore key: images}
\end{itemize}

\subsection{Prompt formatting}
\texttt{math.jinja} enforces a response format where the final answer is placed in \texttt{\textbackslash boxed\{\}} and the model optionally emits a \texttt{<think>} block. This allows the reward function to reliably extract the final answer.

\subsection{Multimodal inputs}
When the example includes images:
\begin{itemize}
  \item The prompt string is split by \texttt{<image>} and converted into a chat message list with interleaved \texttt{\{type: image\}} and \texttt{\{type: text\}} items.
  \item Image paths/bytes are not converted to tensors in the dataset; they are stored in \texttt{multi	extunderscore modal	extunderscore data}.
  \item The processor later converts images into tensors during the training worker step (see \texttt{rlvr/verl/workers/fsdp	extunderscore workers.py}).
\end{itemize}

\section{Reward Function: Tesserae}
\begin{itemize}
  \item Reward entry point: \texttt{rlvr/examples/reward\_function/reward\_tesserae.py:compute\_score}.
  \item Reward type: \texttt{batch} (configured in \texttt{rlvr/examples/config.yaml}).
  \item Reward manager: \texttt{rlvr/verl/workers/reward/function.py}.
\end{itemize}

\subsection{Format reward}
A format reward (optional) is computed by checking whether the response matches a regex containing \texttt{\textbackslash boxed\{...\}}. In the example config, \texttt{format\_weight} is set to \texttt{0.0}, so the explicit format reward is not mixed in. However, the accuracy computation is gated on format: if the response does not contain \texttt{\textbackslash boxed\{...\}}, accuracy is forced to \(0\).

\subsection{Grid accuracy for Tesserae}
The reward function tries to parse the boxed content into a 2D binary grid (0/1). It accepts:
\begin{itemize}
  \item A Python-like list string (e.g. \texttt{[[1,0],[0,1]]}).
  \item A literal list-of-lists.
  \item A NumPy 2D array.
\end{itemize}
If parsing fails or the shape mismatches the ground truth, reward is \(0.0\).

When valid, the grid accuracy score is:
\begin{equation}
  r = \frac{\text{matches}}{\text{total}}
\end{equation}
This yields a range of \([0, 1]\), where \(0.5\) corresponds to random 50/50 guessing.

\subsection{Fallback for non-grid data}
If ground truth is not a binary grid, the reward falls back to \texttt{mathruler.grader.grade\_answer} on the boxed content.

\section{Token-level Reward Tensor}
The reward manager returns a \emph{token-level} tensor, but only the last response token receives the scalar score:
\begin{itemize}
  \item Shape: \texttt{(batch, response\_length)}.
  \item All entries are \(0\) except the last response token.
\end{itemize}

Example for response length 5 with score \(0.6\):
\[ [0, 0, 0, 0, 0.6] \]

Sequence-level reward is the sum across tokens, which equals the scalar score.

\section{GRPO Advantage Construction}
\begin{itemize}
  \item Advantage estimator: \texttt{grpo} (\texttt{rlvr/verl/trainer/core\_algos.py}).
  \item Requires \texttt{rollout.n > 1} so each prompt has multiple sampled responses.
\end{itemize}

Let \(R_i\) be the scalar reward for response \(i\), obtained by summing the token-level reward tensor. Responses are grouped by a shared \texttt{uid} (one uid per prompt). GRPO normalizes within each group:
\begin{equation}
  A_i = \frac{R_i - \mu_{g}}{\sigma_{g} + \epsilon}
\end{equation}
where \(g\) is the prompt group, and \(\mu_g\), \(\sigma_g\) are the mean and std of rewards in that group.

The normalized scalar is then broadcast to response tokens:
\begin{equation}
  \text{returns} = A_i \cdot \text{response\_mask}
\end{equation}
Thus, \emph{all response tokens share the same advantage} for a given response, and prompt tokens are masked out.

\section{Forward and Backward Pass (Multimodal)}
This section describes what happens when a prompt has both text and images.

\subsection{Step 1: Prepare multimodal inputs}
\begin{itemize}
  \item In \texttt{fsdp\_workers.py}, \texttt{\_process\_multi\_modal\_inputs} converts images into tensors using the model processor (e.g., Qwen-VL image processor).
  \item These tensors are stored in \texttt{multi\_modal\_inputs} and passed into the model forward.
\end{itemize}

\subsection{Step 2: Model forward (logits)}
\begin{itemize}
  \item The actor forward call uses \texttt{input\_ids}, \texttt{position\_ids}, and \texttt{multi\_modal\_inputs}.
  \item The vision tower produces image embeddings, which are fused with the text tokens (via the model architecture).
  \item Logits are computed for the full sequence, then sliced to response positions.
\end{itemize}

\subsection{Step 3: Policy loss}
\begin{itemize}
  \item PPO/GRPO loss is computed only on response tokens via \texttt{response\_mask}.
  \item The same scalar advantage for a response is applied to all its response tokens.
\end{itemize}

\subsection{Step 4: Backpropagation}
Even though the loss is only on response tokens, gradients flow through the entire computation graph:
\begin{itemize}
  \item Response logits depend on prompt context (text + image embeddings).
  \item Therefore, gradients flow back through the fusion layers into the vision tower.
  \item If the vision tower is trainable, its parameters receive updates.
\end{itemize}

\section{Vision Tower Updates}
\begin{itemize}
  \item Default: \texttt{freeze\_vision\_tower = false} (in \texttt{rlvr/examples/config.yaml}), so vision parameters are updated.
  \item If \texttt{freeze\_vision\_tower = true}, FSDP sets \texttt{requires\_grad\_(False)} on the vision module (see \texttt{rlvr/verl/workers/fsdp\_workers.py}).
\end{itemize}

\section{KL Handling}
There are two modes:
\begin{itemize}
  \item \textbf{use\_kl\_loss = false:} KL penalty is subtracted directly from token-level rewards.
  \item \textbf{use\_kl\_loss = true:} KL term is added to the policy loss (default in the example config).
\end{itemize}

\section{End-to-End Summary}
\begin{enumerate}
  \item Load example (text + image) and build prompt.
  \item Generate responses (\texttt{rollout.n} per prompt).
  \item Compute scalar reward per response using \texttt{reward\_tesserae}.
  \item Place scalar reward on last response token; sum to get \(R_i\).
  \item GRPO normalizes rewards within each prompt group to get advantages.
  \item Forward pass with multimodal inputs; compute response log-probs.
  \item Compute policy loss on response tokens; backprop through entire model.
  \item Update all trainable params, including vision tower unless frozen.
\end{enumerate}

\end{document}
